\documentclass[11pt]{article}
\usepackage[margin=1in]{geometry}
\usepackage{amsmath,amssymb,amsthm}
\usepackage[T1]{fontenc}
\usepackage[utf8]{inputenc}
\usepackage{graphicx}
\usepackage{hyperref}

\newtheorem{theorem}{Theorem}

\newcommand{\esssup}{\operatorname*{ess\,sup}}
\newcommand{\supp}{\operatorname{supp}}
\newcommand{\norm}[1]{\left\lVert #1 \right\rVert}
\newcommand{\tX}{\widetilde X}
\newcommand{\tY}{\widetilde Y}

\title{A compact-multiplier obstruction for Tandori function spaces}
\author{Run \texttt{fa\_banach\_001}}
\date{2026-06-14}

\begin{document}
\maketitle

\section*{Source Problem}

This packet concerns Problem 7.5 of Kiwerski--Tomaszewski
\cite{KT}, page 51 of the source PDF:

\begin{center}
\includegraphics[width=\textwidth]{figures/open_problem_crop.png}
\end{center}

The problem asks what can be said about weakly compact and compact
multipliers between abstract Tandori function spaces.

\section*{Proof Intuition}

The Tandori norm sees the right endpoint of a set rather than its exact
location. If a measurable set \(A\subset (0,b)\) has essential supremum
\(b\), then the least decreasing majorant of \(\chi_A\) is
\(\chi_{(0,b)}\). Thus one can split any positive-measure level set of a
nonzero symbol into infinitely many disjoint pieces with the same essential
supremum. Their characteristic functions form a bounded sequence in
\(\widetilde X\), while the images under multiplication by the symbol remain
uniformly separated in \(\widetilde Y\). This rules out compactness.

\section*{Theorem}

Let \(I=(0,1)\) or \(I=(0,\infty)\), equipped with Lebesgue measure. Let
\(X\) and \(Y\) be Köthe function spaces on \(I\), in the convention of
\cite{KT}, and define the corresponding Tandori function spaces by
\[
  \widetilde X=\{f:\widetilde f\in X\},\qquad
  \norm{f}_{\widetilde X}=\norm{\widetilde f}_{X},
\]
where
\[
  \widetilde f(x)=\esssup_{t\ge x}|f(t)|.
\]
If \(M_\lambda:\widetilde X\to \widetilde Y\) is a compact multiplication
operator, then \(\lambda=0\) almost everywhere.

\section*{Proof}

Assume, toward a contradiction, that \(\lambda\ne 0\) on a set of positive
measure. Then for some \(\varepsilon>0\) the set
\[
  E=\{t\in I: |\lambda(t)|>\varepsilon\}
\]
has positive measure. Since \(I\) is \(\sigma\)-finite, there is a finite
number \(N\) such that \(E\cap(0,N)\) has positive measure; in the case
\(I=(0,1)\) take \(N<1\) after replacing \(E\) by a positive-measure
intersection with a smaller interval if necessary. Put
\[
  B=E\cap(0,N),\qquad b=\operatorname*{ess\,sup} B .
\]
Then \(0<b<\infty\), and \(B\cap(b-\delta,b)\) has positive measure for
every \(\delta>0\).

Choose infinitely many pairwise disjoint measurable subsets
\(B_1,B_2,\ldots\subset B\), each with essential supremum \(b\). For
completeness, here is one construction. Pick infinitely many disjoint
positive-measure pieces in annuli approaching \(b\), and partition the
corresponding annulus indices into countably many infinite sets; the union of
the pieces in the \(n\)-th infinite block is \(B_n\). Non-atomicity lets us
split the positive-measure annuli as needed.

For each \(n\), since \(\operatorname*{ess\,sup} B_n=b\), the least decreasing
majorant of \(\chi_{B_n}\) is exactly \(\chi_{(0,b)}\). Set
\[
  u_n=\frac{\chi_{B_n}}{\norm{\chi_{(0,b)}}_X}.
\]
The source convention for Köthe spaces ensures \(\chi_{(0,b)}\in X\), so
\(\norm{u_n}_{\widetilde X}=1\).

If \(n\ne m\), then \(B_n\cup B_m\subset E\) is still a set with essential
supremum \(b\). Therefore
\[
  \widetilde{\lambda(u_n-u_m)}(x)
  \ge \frac{\varepsilon}{\norm{\chi_{(0,b)}}_X}\chi_{(0,b)}(x)
  \quad\text{for a.e. }x\in I.
\]
By the ideal property of \(Y\),
\[
  \norm{M_\lambda u_n-M_\lambda u_m}_{\widetilde Y}
  \ge
  \varepsilon\,
  \frac{\norm{\chi_{(0,b)}}_Y}{\norm{\chi_{(0,b)}}_X}
  >0 .
\]
Thus the bounded sequence \((u_n)\) in \(\widetilde X\) has images under
\(M_\lambda\) with no Cauchy subsequence in \(\widetilde Y\). This contradicts
compactness of \(M_\lambda\). Hence \(\lambda=0\) almost everywhere.

\section*{Verification And Scope}

The proof is purely analytic and uses only the source paper's definition of
Köthe function spaces and Tandori function spaces, plus the non-atomicity of
the underlying interval. It answers the compact-multiplier part of Problem
7.5 in the standard function-space setting \(I=(0,1)\) or \(I=(0,\infty)\).
It does not settle the weakly compact multiplier part of the problem.

\section*{Novelty Check}

Local run indexes contained no packet or attempt for arXiv:2212.06723 or for
the Tandori compact-multiplier problem. Bounded web searches on 2026-06-14 for
exact and close phrases such as \texttt{"compact multipliers" "Tandori"
function spaces}, \texttt{"Tandori function space" "compact multipliers"}, and
\texttt{"weakly compact multipliers" "Tandori"} found the source paper but no
later solution of Problem 7.5.

\begin{thebibliography}{9}

\bibitem{KT}
T. Kiwerski and J. Tomaszewski,
\emph{Essential norms of pointwise multipliers in the non-algebraic setting},
arXiv:2212.06723.

\end{thebibliography}

\end{document}
